% =====================================================================
%  BÁO CÁO GỌN 2–3 TRANG — LAB 01: SQL TRÊN CLOUD SHELL VỚI SQLITE
%  Introduction to Data Science — PTIT
%
%  Đây là bản NỘP CHÍNH (đúng giới hạn 2–3 trang của đề). Bản đầy đủ có phụ lục
%  nằm ở lab01_report.tex, dùng khi không bị giới hạn số trang.
%
%  Biên dịch:
%      bash compile.sh 2p          (hoặc: pdflatex lab01_report_short.tex, 2 lần)
%  Bố cục 2 cột, 10pt, lề 1.6cm, bảng \footnotesize — đã tính để vừa 3 trang.
%  Nếu máy bạn ra 4 trang: xem mục "Co giãn độ dài" trong report/README.md.
% =====================================================================
\documentclass[10pt,a4paper,twocolumn]{article}

\usepackage[utf8]{inputenc}
\usepackage[T5]{fontenc}
\usepackage[vietnamese]{babel}

\usepackage[a4paper,top=1.5cm,bottom=1.5cm,left=1.6cm,right=1.6cm,columnsep=0.65cm]{geometry}
\usepackage{amsmath}
\usepackage{booktabs}
\usepackage{tabularx}
\usepackage{array}
\usepackage{enumitem}
\usepackage{float}
\usepackage{caption}
\usepackage{listings}
\usepackage{xcolor}
\usepackage{titlesec}
\usepackage[hidelinks,unicode]{hyperref}

% ---------- Nén khoảng cách để vừa 3 trang ----------
\setlength{\parindent}{0pt}
\setlength{\parskip}{2.5pt}
\renewcommand{\arraystretch}{1.08}
\captionsetup{font=scriptsize,labelfont=bf,skip=3pt}
\titlespacing*{\section}{0pt}{6pt}{3pt}
\titlespacing*{\subsection}{0pt}{4pt}{2pt}
\titleformat{\section}{\normalfont\bfseries\large}{\thesection.}{0.4em}{}
\titleformat{\subsection}{\normalfont\bfseries\normalsize}{\thesubsection}{0.4em}{}
\setlist[itemize]{leftmargin=1.1em,itemsep=1pt,topsep=2pt,parsep=0pt}
\setlist[enumerate]{leftmargin=1.3em,itemsep=1pt,topsep=2pt,parsep=0pt}

\renewcommand{\figurename}{Hình}
\renewcommand{\tablename}{Bảng}
\newcommand{\code}[1]{\texttt{\small #1}}

\definecolor{codebg}{gray}{0.95}
\lstset{
  basicstyle=\ttfamily\scriptsize,
  backgroundcolor=\color{codebg},
  frame=single,
  rulecolor=\color{gray!60},
  keywordstyle=\color{blue!60!black}\bfseries,
  commentstyle=\color{green!35!black}\itshape,
  showstringspaces=false,
  breaklines=true,
  columns=fullflexible,
  keepspaces=true,
  aboveskip=3pt,
  belowskip=3pt,
  xleftmargin=1pt,
  captionpos=b
}

\begin{document}

% =====================================================================
%  TIÊU ĐỀ (trải hết chiều ngang, không chia cột)
% =====================================================================
\twocolumn[{
\begin{center}
{\footnotesize TRƯỜNG ĐẠI HỌC BƯU CHÍNH VIỄN THÔNG (PTIT) \quad$\cdot$\quad KHOA CÔNG NGHỆ THÔNG TIN\par}
\vspace{3pt}
{\Large\bfseries BÁO CÁO THỰC HÀNH LAB 01\par}
\vspace{1pt}
{\large SQL trên Google Cloud Shell với SQLite\par}
\vspace{2pt}
{\footnotesize Môn: Introduction to Data Science \quad$\cdot$\quad
GV: \dotfill\ \ SV: \dotfill\ \ MSSV: \dotfill\ \ Lớp: \dotfill\par}
\vspace{2pt}
{\color{black!30}\rule{\linewidth}{0.4pt}}
\end{center}

\noindent\textbf{Tóm tắt.} Báo cáo trình bày việc xây dựng cơ sở dữ liệu quan hệ bán hàng trực
tuyến bốn bảng (\code{customers}, \code{products}, \code{orders}, \code{order\_items}) trên
SQLite trong Google Cloud Shell, thực hiện CRUD, \code{JOIN}, \code{GROUP BY}/\code{HAVING},
subquery, window function, rồi dùng index và \code{EXPLAIN QUERY PLAN} để chứng minh ảnh hưởng
của index tới kế hoạch thực thi. Cả 21 yêu cầu (S01--S21) đã cài đặt và chạy kiểm chứng thực tế:
0 lỗi, dữ liệu luôn thoả khoá ngoại, kế hoạch chuyển từ \code{SCAN orders} sang
\code{SEARCH orders USING COVERING INDEX} sau khi tạo index.

\vspace{3pt}
}]

% =====================================================================
\section{Mục tiêu và môi trường}
% =====================================================================

Bài thực hành mô phỏng cửa hàng trực tuyến nhỏ với dữ liệu chuẩn hoá thành bốn bảng, nhằm trả lời
các câu hỏi kinh doanh cơ bản (đơn giá trị cao nhất, khách chi tiêu nhiều nhất, sản phẩm bán
chạy, doanh thu theo tháng) và rèn kỹ năng đọc kế hoạch truy vấn. Môi trường thực hiện hoàn toàn
miễn phí, không bật Billing.

\begin{table}[H]
\centering
\caption{Môi trường thực hành}
\label{tab:env}
\scriptsize
\begin{tabularx}{\linewidth}{@{}l X@{}}
\toprule
Nền tảng & Google Cloud Shell (chế độ \code{\$HOME} bền vững) \\
Python & 3.11.2, chỉ dùng thư viện chuẩn \code{sqlite3} \\
CSDL & SQLite 3.40.1 (window function yêu cầu $\geq$ 3.25) \\
Công cụ & \code{sql\_console.py}, \code{build\_db.py}, \code{generate\_data.py}, \code{run\_sql.py} \\
\bottomrule
\end{tabularx}
\end{table}

\textbf{Phân chia nội dung.} Đề có hai khối yêu cầu với điều kiện dữ liệu khác nhau nên bài làm
tách hai tệp đáp án riêng: \textbf{bắt buộc} S01--S15 trong \code{sql/answers.sql} chạy trên
\code{data/sales.db} (dữ liệu mẫu của đề); \textbf{bổ sung/về nhà} S16--S21 trong
\code{sql/answers\_homework.sql} chạy trên \code{data/sales\_extended.db} (dữ liệu mở rộng);
phần chứng minh index ở \code{sql/explain.sql}. Cách tách này giúp người chấm thấy ngay đâu là
yêu cầu tối thiểu, đâu là phần mở rộng.

% =====================================================================
\section{Mô hình dữ liệu}
% =====================================================================

Bốn bảng với ba quan hệ một--nhiều: \code{customers} (1--N) \code{orders},
\code{orders} (1--N) \code{order\_items}, \code{products} (1--N) \code{order\_items}.

\begin{lstlisting}[language=SQL,caption={Trích \code{sql/schema.sql}}]
CREATE TABLE customers (
  customer_id INTEGER PRIMARY KEY,
  name TEXT NOT NULL, city TEXT);

CREATE TABLE products (
  product_id INTEGER PRIMARY KEY,
  name TEXT NOT NULL, category TEXT,
  price REAL CHECK(price >= 0));

CREATE TABLE orders (
  order_id INTEGER PRIMARY KEY,
  customer_id INTEGER NOT NULL, order_date TEXT,
  FOREIGN KEY(customer_id) REFERENCES customers(customer_id));

CREATE TABLE order_items (
  order_id INTEGER, product_id INTEGER,
  quantity INTEGER CHECK(quantity > 0),
  PRIMARY KEY(order_id, product_id),
  FOREIGN KEY(order_id) REFERENCES orders(order_id),
  FOREIGN KEY(product_id) REFERENCES products(product_id));
\end{lstlisting}

Ba quyết định thiết kế đáng chú ý: (i) khoá chính tổ hợp \code{(order\_id, product\_id)} ngăn một
sản phẩm xuất hiện hai lần trong cùng đơn --- muốn mua thêm thì tăng \code{quantity};
(ii) \code{CHECK} chặn giá âm và số lượng $\leq 0$ ngay tại tầng CSDL;
(iii) SQLite \emph{không} kiểm tra khoá ngoại mặc định, nên mọi kết nối trong bài đều chạy
\code{PRAGMA foreign\_keys = ON} ngay sau khi mở CSDL --- thiếu bước này thì ràng buộc chỉ còn là
chú thích.

% =====================================================================
\section{Kết quả phần bắt buộc (S01--S15)}
% =====================================================================

Chạy \code{python3 run\_sql.py sql/answers.sql} thực thi 18 câu lệnh (gồm câu kiểm tra trung
gian) với \textbf{0 lỗi}; chạy lần thứ hai vẫn 0 lỗi và cho kết quả giống hệt. Dữ liệu mẫu sau
khi nạp khớp đúng kỳ vọng của đề: \code{customers}\,5, \code{products}\,5, \code{orders}\,6,
\code{order\_items}\,8; \code{PRAGMA foreign\_key\_check} trả về rỗng.

\begin{table}[H]
\centering
\caption{Tổng hợp 15 câu bắt buộc (số dòng là kết quả thực tế)}
\label{tab:s01s15}
\scriptsize
\begin{tabularx}{\linewidth}{@{}l X r@{}}
\toprule
\textbf{Mã} & \textbf{Kỹ thuật} & \textbf{Dòng} \\
\midrule
S01 & \code{SELECT} + \code{ORDER BY} & 5 \\
S02 & \code{WHERE price > 1000000} & 2 \\
S03 & \code{ORDER BY price DESC} & 5 \\
S04 & \code{INSERT ... ON CONFLICT DO NOTHING} & 1 \\
S05 & \code{UPDATE ... SET price = price*1.1} & 1 \\
S06 & \code{DELETE} + \code{NOT EXISTS} & 1 \\
S07 & \code{INNER JOIN} hai bảng & 6 \\
S08 & \code{JOIN} ba bảng, \code{GROUP BY}, \code{SUM} & 6 \\
S09 & \code{LEFT JOIN} + \code{GROUP BY} & 5 \\
S10 & \code{HAVING SUM(...) > 5000000} & 3 \\
S11 & subquery vô hướng \code{> (SELECT AVG(price))} & 1 \\
S12 & \code{HAVING COUNT(order\_id) >= 2} & 1 \\
S13 & \code{ORDER BY} + \code{LIMIT 3} & 3 \\
S14 & \code{CREATE INDEX} trên \code{orders(customer\_id)} & --- \\
S15 & \code{EXPLAIN QUERY PLAN} & 1 \\
\bottomrule
\end{tabularx}
\end{table}

\textbf{CRUD lặp an toàn.} Ba câu có ghi (S04--S06) được viết để file chạy lại không sai dữ liệu:
\code{INSERT ... ON CONFLICT(customer\_id) DO NOTHING}; \code{UPDATE} kèm điều kiện
\code{AND price = 250000} nên không cộng dồn 10\% ở lần chạy sau (lần hai trả về
\code{Rows affected: 0}); \code{DELETE} kèm \code{NOT EXISTS} đúng ngữ nghĩa "khách vừa thêm chưa
có đơn" và đồng thời bảo vệ khoá ngoại. Kiểm chứng sau hai lần chạy: giá Mouse vẫn
\textbf{275.000}, số khách trở về \textbf{5}.

\begin{table}[H]
\centering
\caption{S08 (tổng tiền từng đơn) và S09 (doanh thu từng khách)}
\label{tab:s08s09}
\scriptsize
\begin{tabular}{@{}r r@{\hspace{6pt}}l r@{}}
\toprule
\textbf{order} & \textbf{total} & \textbf{Khách} & \textbf{revenue} \\
\midrule
1001 & 15.550.000 & An (1)   & 17.150.000 \\
1004 & 15.000.000 & Chi (3)  & 15.000.000 \\
1006 &  7.000.000 & Ha (5)   &  7.000.000 \\
1002 &  4.150.000 & Binh (2) &  4.150.000 \\
1003 &  1.600.000 & Dung (4) &    825.000 \\
1005 &    825.000 &          &           \\
\midrule
\multicolumn{1}{@{}r}{Tổng} & 44.125.000 & & \\
\bottomrule
\end{tabular}
\end{table}

Kiểm tra chéo: đơn 1001 gồm 1 Laptop (15.000.000) và 2 Mouse (2\,$\times$\,275.000) =
15.550.000, khớp kết quả máy. Con số này cho thấy \textbf{thứ tự thực hiện có ảnh hưởng}: vì S05
đã tăng giá Mouse trước khi chạy S08, tổng dùng giá 275.000 chứ không phải 250.000.

\textbf{Nhóm truy vấn nâng cao.} S10 lọc theo kết quả tổng hợp nên phải dùng \code{HAVING}
(không phải \code{WHERE}) --- đúng 3 khách An, Chi, Ha vượt 5 triệu. S11 so sánh với giá trung
bình 4.045.000 (sau khi S05 đổi giá) nên chỉ Laptop vượt trung bình; subquery không tương quan
nên SQLite tính một lần. S12 dùng \code{HAVING COUNT(o.order\_id) >= 2}, kết quả đúng một khách:
An với 2 đơn. S13 lấy top 3 theo doanh thu: Laptop 30.000.000, Monitor 10.500.000, Headphones
1.600.000 --- nhưng \code{LIMIT 3} không xử lý đồng hạng, muốn đúng phải dùng \code{RANK()} như
S19.

% =====================================================================
\section{Index và kế hoạch thực thi (S14--S15)}
% =====================================================================

Vì S15 yêu cầu so sánh \emph{trước} và \emph{sau} index, \code{sql/explain.sql} tự đưa CSDL về
trạng thái cần thiết: \code{DROP INDEX IF EXISTS} $\rightarrow$ \code{EXPLAIN} $\rightarrow$
\code{CREATE INDEX IF NOT EXISTS} $\rightarrow$ \code{EXPLAIN}, nhờ vậy file chạy đúng trên bất kỳ
bản sao CSDL nào.

\begin{table}[H]
\centering
\caption{Kế hoạch thực thi trước/sau index \code{idx\_orders\_customer}}
\label{tab:explain}
\scriptsize
\begin{tabularx}{\linewidth}{@{}l X@{}}
\toprule
Trước & \code{SCAN orders} \\
Sau & \code{SEARCH orders USING COVERING INDEX idx\_orders\_customer (customer\_id=?)} \\
\bottomrule
\end{tabularx}
\end{table}

\code{SCAN} đọc tuần tự toàn bảng rồi mới lọc --- chi phí tăng tuyến tính theo số dòng.
\code{SEARCH} đi xuống cây B-tree của index, chi phí $\sim\log n$ cộng số dòng khớp.
\code{COVERING INDEX} là chi tiết đáng chú ý: câu chỉ cần \code{order\_id}, mà index trên
\code{customer\_id} của bảng có \code{INTEGER PRIMARY KEY} đã chứa sẵn rowid (chính là
\code{order\_id}), nên SQLite trả lời hoàn toàn từ index, không quay lại bảng; nếu lấy thêm
\code{order\_date} (câu S20) thì kế hoạch là \code{SEARCH ... USING INDEX}.

\textbf{Không kết luận bằng thời gian chạy.} Với 6 đơn, chênh lệch \code{SCAN}/\code{SEARCH} nằm
dưới ngưỡng đo của đồng hồ hệ thống --- so sánh wall-clock sẽ kết luận sai là index "không có tác
dụng". Cách đúng là đọc \emph{kế hoạch} (kiểu truy cập, số bước, có phải sort thêm hay không).

% =====================================================================
\section{Kết quả phần bổ sung (S16--S21)}
% =====================================================================

\textbf{Dữ liệu mở rộng.} \code{generate\_data.py} dùng \code{random.seed(2026)} để tái lập được
và ghi vào \code{data/sales\_extended.db} (không ghi đè \code{sales.db}): \textbf{100} khách,
\textbf{30} sản phẩm, \textbf{500} đơn, \textbf{1.500} dòng đơn (mỗi đơn 1--5 sản phẩm), vượt mức
tối thiểu của đề (100/30/500/1.000), \code{foreign\_key\_check} \code{ok}. Kịch bản chủ động để
$\sim$8\% khách không phát sinh đơn và 4 sản phẩm cuối không xuất hiện trong đơn nào, nhằm bảo
đảm S17, S18 có kết quả.

\begin{itemize}
  \item \textbf{S16} doanh thu theo tháng (\code{strftime('\%Y-\%m', ...)}) cho 13 nhóm tháng
        (10/2025--10/2026); cao nhất 2025-10 với 1.142.485.000\,đ / 53 đơn. Tháng cuối chỉ có 7
        đơn --- hiện tượng \emph{tháng dở dang}, cần nêu rõ kỳ dữ liệu khi báo cáo.
  \item \textbf{S17} khách chưa có đơn: \code{LEFT JOIN ... WHERE o.order\_id IS NULL} → 8 khách
        (id 37, 38, 43, 56, 83, 89, 92, 94).
  \item \textbf{S18} sản phẩm chưa từng bán: \code{NOT EXISTS} → 4 sản phẩm (Graphics Card,
        RAM 16GB, CPU Cooler, Power Supply). Cả S17, S18 đều là bài toán "không có bản ghi tương
        ứng"; \code{NOT EXISTS} được ưu tiên vì diễn đạt trực tiếp ngữ nghĩa.
  \item \textbf{S19} xếp hạng bằng window function: top 3 là Ha 314.565.000, Ho Mai 281.910.000,
        Hoang Chi 267.400.000. Dữ liệu này không có đồng hạng nên \code{RANK} và \code{DENSE\_RANK}
        trùng nhau; khi có hai khách cùng hạng 2, \code{RANK} nhảy sang 4 còn \code{DENSE\_RANK}
        sang 3 --- báo cáo "top N" nên dùng \code{DENSE\_RANK}.
  \item \textbf{S20} sau khi mở rộng dữ liệu: chưa index là \code{SCAN orders}, có index là
        \code{SEARCH orders USING INDEX idx\_orders\_customer} kèm
        \code{USE TEMP B-TREE FOR ORDER BY} (do \code{ORDER BY order\_date} trên cột không nằm
        trong index). Nếu truy vấn chạy thường xuyên, index tổ hợp
        \code{(customer\_id, order\_date)} sẽ loại bỏ bước sort.
  \item \textbf{S21} so sánh bảng tạm: SQL Server \code{\#t} là bảng tạm cục bộ theo session,
        \code{\#\#t} là bảng tạm toàn cục dùng chung giữa các session, cả hai sống trong
        \code{tempdb} (CSDL hệ thống, cũng chứa biến bảng và sort/hash spill). SQLite chỉ có
        \code{CREATE TEMP TABLE} gắn với \emph{kết nối}, nằm trong schema \code{temp} (RAM hoặc
        tệp tạm tuỳ \code{PRAGMA temp\_store}) và \textbf{không} có bảng tạm toàn cục; bảng
        \code{temp} được tìm trước \code{main} nên trùng tên sẽ che bảng thường. Đã minh hoạ bằng
        \code{TEMP TABLE} tổng hợp doanh thu tháng rồi \code{DROP TABLE temp.}.
\end{itemize}

% =====================================================================
\section{Kết luận và bài học}
% =====================================================================

Cả 21 yêu cầu đã hoàn thành: CSDL bốn bảng đúng thiết kế với khoá chính/ngoại và \code{CHECK};
dữ liệu mẫu 5/5/6/8; S01--S15 và S16--S21 chạy không lỗi và lặp an toàn; index được chứng minh
bằng kế hoạch thực thi. Bốn bài học: (i) khoá ngoại SQLite là tuỳ chọn --- phải bật pragma;
(ii) script có ghi phải lặp an toàn (\code{ON CONFLICT}, điều kiện chặt trong \code{UPDATE},
\code{NOT EXISTS} trong \code{DELETE}); (iii) index không miễn phí --- chỉ tạo cho cột thường
lọc/nối; (iv) đọc plan thay vì đo thời gian, đặc biệt chú ý \code{COVERING INDEX} và
\code{USE TEMP B-TREE FOR ORDER BY} vì chúng chỉ ra chi phí ẩn.

\begin{lstlisting}[language=bash,caption={Lệnh tái lập toàn bộ kết quả}]
python3 build_db.py --force      # 5 / 5 / 6 / 8
python3 run_sql.py sql/answers.sql
python3 generate_data.py --force # 100 / 30 / 500 / 1500
python3 run_sql.py sql/answers_homework.sql \
        --db data/sales_extended.db
python3 run_sql.py sql/explain.sql  # SCAN -> SEARCH
\end{lstlisting}

\end{document}
